\clearpage
\section{Modeling biological processes: variables, rates, and time evolution}
\label{sec:rates}
% EDITORIAL ADDITION: navigation or conceptual bridge
A model starts with a question and a choice of variables.
For growth, a single mass variable may suffice; for binding, concentrations of free and bound molecules are useful.
The aim is to choose a description that retains the processes needed to answer the question.
The broader problem of finding useful representations is revisited in section~\ref{sec:representations}.
% Source: day2.tex:15; unit d2.1.0
\subsection{The theory of the mean}
\label{sec:mean}


% ============================================================

A useful way of constructing a dynamical model is to keep track of the mean value of one or more variables of interest.

The basic protocol is:

\begin{enumerate}
    \item Identify the variables of interest.
    \item Enumerate all the ways in which those variables can increase or decrease, and assign rates to those processes.
    \item Write an update rule describing the change during a small time interval $\Delta t$.
    \item Take the continuum limit and integrate the resulting differential equation to obtain the time evolution.
\end{enumerate}


% ------------------------------------------------------------

% Source: day1.tex:69; unit d1.3
\subsection{Autocatalysis and exponential cellular growth}
\label{sec:growth}


% ====================

A useful coarse-grained picture of the cell is as a system whose dry mass is dominated by proteins and whose growth is therefore closely related to protein production \cite{milo2013total}.
In \textit{E.~coli}, measured protein-to-mRNA abundance ratios span approximately $10^2$ - $10^4$ across genes, making
\[
    \frac{\text{number of protein molecules}}{\text{number of mRNA molecules}} \sim 10^3
\]
a useful order-of-magnitude scale \cite{taniguchi2010quantifying}.

Ribosomes synthesize proteins, including the ribosomal proteins required to construct new ribosomes.
Thus the translational machinery participates in its own production.
At a coarse-grained level, this creates an \emph{autocatalytic} process: the machinery responsible for growth produces more of the machinery responsible for growth \cite{roy2021unifying}.
This provides an intuitive route to exponential cellular growth.

% ====================
\subsubsection*{Simple growth law}
% ====================

Let $M(t)$ denote the cell mass at time $t$.

Assume that the rate at which new mass is produced is proportional to the mass already present. 
Over a short time interval $\Dt$, 
\[
    M(t+\Dt) = M(t) + kM(t)\Dt,
\] 
where $k$ is the fractional growth rate. 
Rearranging, 
\[
    \frac{M(t+\Dt)-M(t)}{\Dt} = kM(t).
\] 
Taking the limit $\Dt\rightarrow 0$ gives 
\[
    \frac{\dd M}{\dd t} = kM.
\] 
The solution is 
\[
    M(t) = M(0)e^{kt}.
\] 
Thus a constant fractional growth rate produces exponential growth.

This model assumes that

\begin{itemize}
    \item the cell can be described by a single coarse-grained variable $M(t)$;
    \item the amount of active protein-production machinery is
    approximately proportional to cell mass;
    \item the protein-production rate per unit mass is approximately
    constant;
    \item nutrients and other resources are not limiting over the time interval considered;
    \item degradation and other losses are either negligible or already incorporated into the effective rate constant $k$;
    \item fluctuations in individual molecular events can be neglected, so that a deterministic continuum description is appropriate.
\end{itemize}


% ====================

% Source: day1.tex:132; unit d1.4
\subsection{Biological turnover}
\label{sec:turnover}


% ====================

Biological systems continually replace components such as red blood cells, photoreceptor components, proteins, RNA molecules, and cells in renewing tissues \cite{reddien2024purpose}.
Human red blood cells provide a classic example: isotope-labeling experiments established that they have a finite lifespan and are continuously replaced \cite{shemin1946life}.
A simple birth-death model illustrates how turnover can produce a steady state even while individual components are continuously produced and destroyed.
The same general modeling framework applies: identify the states of interest, enumerate the processes that move material between those states, assign rates, and construct the corresponding dynamical equations.

% Source: day2.tex:265; unit d2.2
\subsection{Example: phytoplankton resource competition}
\label{sec:plankton}
A biological example of constructing dynamical rules from rates is the uptake of limiting resources by phytoplankton.
David Tilman studied competition between freshwater phytoplankton species under different supplies of limiting resources.
The experiments compared predictions from quantitative resource-competition models with the observed abundances of competing algal species \cite{tilman1977resource}.

% Source: day1.tex:147; unit d1.5
\subsection{Case study: mRNA production, degradation, and steady state}
\label{sec:mrna-mean}


% ====================

Let $m(t)$ denote the mean number of mRNA molecules per cell at time $t$.

Suppose that
\begin{itemize}
    \item mRNA molecules are produced at a constant rate $r$;
    \item each existing mRNA has a constant probability per unit time $\gamma$ of being degraded.
\end{itemize}

For example, consider a production rate
\[
    r \sim 1 \text{ mRNA per few minutes},
\]
while the characteristic mRNA lifetime is
\[
    \tau = \frac{1}{\gamma}.
\]
Over a small time interval $\Dt$,
\begin{align*}
    m(t+\Dt) &= m(t) + \text{mRNAs produced} - \text{mRNAs degraded}
    \\
    &= m(t) + r\Dt - \gamma m(t)\Dt.
\end{align*}
Therefore,
\[
    \frac{m(t+\Dt)-m(t)}{\Dt} = r-\gamma m(t).
\]
Taking $\Dt \rightarrow 0$,
\begin{equation}
    \frac{\dd m}{\dd t} = r-\gamma m.
    \label{eq:mrna-mean}
\end{equation}
Dimensional consistency determines the units of the parameters.
The left-hand side has units of molecules per time.
Therefore $r$ must have units of molecules per time and $\gamma$ must have units of inverse time, so that $\gamma m$ also has units of molecules per time.
The product $\gamma t$ is dimensionless, as required when it appears in the exponential $e^{-\gamma t}$.

At steady state,
\[
    \frac{\dd m}{\dd t}=0,
\]
 so
\[
    m_{\mathrm{ss}} = \frac{r}{\gamma}.
\]

The full solution is
\[
    m(t) = \frac{r}{\gamma} + \left( m(0)-\frac{r}{\gamma} \right) e^{-\gamma t}.
\]
Thus the system approaches its steady-state mRNA number on the characteristic timescale
\[
    \tau = \frac{1}{\gamma}.
\]

% ====================
The deterministic model assumes that

\begin{itemize}
    \item transcription occurs at a constant average rate $r$;
    \item degradation is first-order: each mRNA has the same constant degradation probability per unit time;
    \item mRNA molecules act independently;
    \item $r$ and $\gamma$ do not vary with time;
    \item discrete molecular fluctuations are ignored.
\end{itemize}

At the level of individual molecules, transcription and degradation are stochastic processes.
Single-cell measurements directly reveal fluctuations in mRNA copy number and burst-like transcription \cite{golding2005real,taniguchi2010quantifying}.
Eqn.~\eqref{eq:mrna-mean} describes the mean behavior.



% ============================================================

% Source: day2.tex:31; unit d2.1.1
\subsection{Case study: ligand-receptor binding}
\label{sec:binding-rates}

% ------------------------------------------------------------

Consider a ligand $L$ that can bind reversibly to a receptor $R$ to form a ligand-receptor complex $LR$.

The variables of interest are therefore the concentrations

\[
    [R], \qquad [L], \qquad [LR].
\]

Here $[R]$ denotes the concentration of free receptors, $[L]$ the concentration of free ligands, and $[LR]$ the concentration of bound ligand-receptor complexes.
% 
% 
% ------------------------------------------------------------
\paragraph{Step 1: identify the variables}
% ------------------------------------------------------------
A schematic representation of the three states is
\begin{figure}[h]
    \centering
    \begin{tikzpicture}[
        ligand/.style={circle, draw, minimum size=0.42cm, inner sep=0pt},
        % reaction/.style={Latex-Latex, thick}
    ]
        \draw[rounded corners=2pt] (-0.55,-0.35) -- (0.55,-0.35) -- (0.55,0.35) -- (0.23,0.35) .. controls (0.18,0.12) and (-0.18,0.12) .. (-0.23,0.35) -- (-0.55,0.35) -- cycle;
        \node at (0,-0.10) {$R$};
        \node at (1.0,0) {$+$};
        \node[ligand] (L) at (2.0,0) {\scriptsize $L$};

        \draw[-Latex, thick] (2.8,0.10) -- (5.2,0.10);
        \draw[-Latex, thick] (5.2,-0.10) -- (2.8,-0.10);
        
        \node[above] at (4.0,0.10) {$k_{\mathrm{on}}$};
        \node[below] at (4.0,-0.10) {$k_{\mathrm{off}}$};

        \begin{scope}[shift={(6.2,0)}]
            \draw[rounded corners=2pt] (-0.55,-0.35) -- (0.55,-0.35) -- (0.55,0.35) -- (0.23,0.35) .. controls (0.18,0.12) and (-0.18,0.12) .. (-0.23,0.35) -- (-0.55,0.35) -- cycle;
            \node at (0,-0.10) {$R$};
            \node[ligand] at (0,0.36) {\scriptsize $L$};
        \end{scope}
        \node at (7.1,0) {$LR$};
    \end{tikzpicture}
    \caption{Reversible binding of a ligand $L$ to a receptor $R$ to form the complex $LR$.
    The same ligand and receptor symbols are used in the unbound and bound states.}
\end{figure}
% 
% 
% ------------------------------------------------------------
\paragraph{Step 2: enumerate the processes and their rates}
% ------------------------------------------------------------
The reaction is
\[
    R + L \xrightleftharpoons[k_{\mathrm{off}}]{k_{\mathrm{on}}} LR.
\]
The forward reaction is ligand binding.
According to the law of mass action, its rate is proportional to both the free ligand concentration and the free receptor concentration,
\[
    \text{binding rate} = k_{\mathrm{on}}[L][R].
\]

The reverse reaction is ligand unbinding.
Its rate is proportional to the concentration of bound complexes,
\[
    \text{unbinding rate} = k_{\mathrm{off}}[LR].
\]
The association rate constant $k_{\mathrm{on}}$ has units of concentration$^{-1}$ time$^{-1}$.
The dissociation rate constant $k_{\mathrm{off}}$ has units of time$^{-1}$.
% 
% 
% ------------------------------------------------------------
\paragraph{Step 3: write the update rule}
% ------------------------------------------------------------
During a short time interval $\Delta t$, the concentration of bound complexes changes according to
\[
    [LR](t+\Delta t) = [LR](t) + k_{\mathrm{on}}\Delta t[L][R] - k_{\mathrm{off}}\Delta t[LR].
\]
The first additional term represents newly formed complexes.
The second represents complexes lost through dissociation.
Subtracting $[LR](t)$ from both sides and dividing by $\Delta t$ gives
\[
    \frac{[LR](t+\Delta t)-[LR](t)}{\Delta t} = k_{\mathrm{on}}[L][R] - k_{\mathrm{off}}[LR].
\]
Taking the limit $\Delta t\rightarrow0$ gives
\[
    \boxed{\frac{\dd[LR]}{\dd t} = k_{\mathrm{on}}[L][R] - k_{\mathrm{off}}[LR].}
\]
% 
% 
% ------------------------------------------------------------
\paragraph{Steady state}
% ------------------------------------------------------------
At steady state, the mean concentration of bound complexes is no longer changing,
\[
    \frac{\dd[LR]}{\dd t} = 0.
\]
Therefore,
\[
    k_{\mathrm{on}}[L][R] - k_{\mathrm{off}}[LR] = 0.
\]
Hence
\[
    k_{\mathrm{on}}[L][R] = k_{\mathrm{off}}[LR].
\]
Rearranging,
\[
    \frac{[L][R]}{[LR]} = \frac{k_{\mathrm{off}}}{k_{\mathrm{on}}}.
\]
We define the dissociation constant
\[
    \boxed{K_d = \frac{k_{\mathrm{off}}}{k_{\mathrm{on}}}.}
\]
Therefore,
\[
    K_d = \frac{[L][R]}{[LR]}.
\]
Equivalently,
\[
    \frac{[LR]}{[R]} = \frac{[L]}{K_d}.
\]
% 
% 
% ------------------------------------------------------------
\paragraph{Probability that a receptor is bound}
% ------------------------------------------------------------
The total receptor concentration is the sum of free and bound receptors,
\[
    [R]_{\mathrm{tot}} = [R] + [LR].
\]
The fraction of receptors that are bound is therefore
\[
    P_{\mathrm{bound}} = \frac{[LR]}{[R]+[LR]}.
\]
Using
\[
    \frac{[LR]}{[R]} = \frac{[L]}{K_d},
\]
we obtain
\[
    P_{\mathrm{bound}} = \frac{[L]/K_d}{1+[L]/K_d}.
\]
Thus
\begin{equation}
    \boxed{P_{\mathrm{bound}}([L]) = \frac{[L]}{K_d+[L]}.}
    \label{eq:binding-occupancy}
\end{equation}
This expression has several useful limits.

When there is no ligand,
\[
    [L]=0 \quad \Rightarrow \quad P_{\mathrm{bound}}=0.
\]
When the ligand concentration is very large compared with $K_d$,
\[
    [L]\gg K_d \quad \Rightarrow \quad P_{\mathrm{bound}}\rightarrow1.
\]
When
\[
    [L]=K_d,
\]
we have
\[
    P_{\mathrm{bound}} = \frac{1}{2}.
\]
Thus $K_d$ is the ligand concentration at which half of the receptors are occupied in this binding model.
A smaller $K_d$ corresponds to stronger binding, since less ligand is required to reach half occupancy.
% 
% 
% ------------------------------------------------------------
\paragraph{Step 4: time evolution}
% ------------------------------------------------------------
To solve explicitly for the time evolution, suppose that the free ligand concentration $[L]$ remains approximately constant.
The total receptor concentration is conserved,
\[
    [R]_{\mathrm{tot}} = [R]+[LR].
\]
Therefore,
\[
    [R] = [R]_{\mathrm{tot}}-[LR].
\]
Substituting this into the dynamical equation gives
\[
    \frac{\dd[LR]}{\dd t} = k_{\mathrm{on}}[L]\left([R]_{\mathrm{tot}}-[LR]\right)-k_{\mathrm{off}}[LR].
\]
Collecting terms gives
\[
    \frac{\dd[LR]}{\dd t} = k_{\mathrm{on}}[L][R]_{\mathrm{tot}}-\left(k_{\mathrm{on}}[L]+k_{\mathrm{off}}\right)[LR].
\]
The equilibrium concentration is
\[
    [LR]_{\mathrm{eq}} = [R]_{\mathrm{tot}}\frac{[L]}{K_d+[L]}.
\]
The full time-dependent solution is
\[
    [LR](t) = [LR]_{\mathrm{eq}} + \left([LR](0)-[LR]_{\mathrm{eq}}\right)e^{-\left(k_{\mathrm{on}}[L]+k_{\mathrm{off}}\right)t}.
\]
The system therefore relaxes exponentially toward equilibrium with characteristic timescale
\[
    \tau = \frac{1}{k_{\mathrm{on}}[L]+k_{\mathrm{off}}}.
\]
% 
% 
% ------------------------------------------------------------
The model assumes that:
\begin{itemize}
    \item the system is well mixed, so spatial variation in ligand and receptor concentrations can be ignored;
    \item each receptor has a single independent ligand-binding site;
    \item binding follows the law of mass action;
    \item the rate constants $k_{\mathrm{on}}$ and $k_{\mathrm{off}}$ are constant;
    \item receptors and ligands are chemically identical within their respective populations;
    \item there is no cooperativity between receptors;
    \item there are no additional receptor states or competing ligands;
    \item for the explicit time-dependent solution above, the free ligand concentration $[L]$ is treated as constant.
\end{itemize}


% ============================================================

% EDITORIAL ADDITION: navigation or conceptual bridge
The same steady occupancy will be recovered by counting microscopic states in section~\ref{sec:binding-counting}.
The rate description additionally gives the relaxation time; an equilibrium occupancy alone does not determine both rate constants.
